% ================================================================
% APPENDICE — REVISIONE DEL METODO DI RENDICONTAZIONE
% ================================================================
\section{Appendice: revisione del metodo di rendicontazione delle ore}
\label{sec:appendice-rendicontazione}

\subsection{Motivazione}

A seguito del colloqio di revisione avvenuto in data 2026-09-15, il
gruppo ha condotto una riunione di riflessione sul concetto di \textbf{ore
produttive} e sul metodo di rendicontazione adottato durante RTB\textsubscript{G} e
PB\textsubscript{G}. Dalla riflessione è emerso che alcune attività effettivamente svolte,
 come raccolta dati per grafici e tabelle, revisione e
sistemazione di documenti, progettazione preliminare prima dello sviluppo,
lettura di un documento; non producendo un artefatto visibile quanto la stesura di un paragrafo o 
la scrittura di codice non venivano rendicontate dalla maggior parte dei membri.

Il gruppo ha quindi effettuato un'analisi retrospettiva, task per task,
individuando le attività non rendicontate per ciascun ruolo e quantificando
il tempo corrispondente, per ruolo e per persona. Questa appendice riporta
i risultati di tale analisi.

Nel corso della riunione, Sofia ha valutato la propria rendicontazione già
in linea con il concetto di ore produttive maturato dal gruppo, e ha
scelto di non modificarla. Gli altri membri, consapevoli di non aver
compreso fino a quel momento quali attività fossero effettivamente da
rendicontare per il proprio ruolo, hanno invece riallineato le ore svolte
alla nuova comprensione. 

Le mancate comprensioni individuate, per ruolo,
sono le seguenti:

\begin{itemize}
    \item \textbf{Amministratore}: non era stato rendicontato il tempo dedicato a 
    raccogliere i dati per i grafici del Piano di Qualifica e per
    le tabelle dei consuntivi del Piano di Progetto, né il tempo di
    revisione e sistemazione di paragrafi ed errori riscontrati nella
    gestione delle Norme di Progetto.
    \item \textbf{Analista}: non era stato considerato il tempo di creazione dei
    grafici dei casi d'uso, né il tempo di rinumerazione manuale dei casi
    d'uso conseguente all'aggiunta o all'eliminazione di uno o più casi
    d'uso.
    \item \textbf{Progettista}: non era stato considerato il tempo per la
    creazione e la modifica dei diagrammi delle classi, dei package, delle
    sequenze e di deployment, dato che anche una singola modifica di
    denominazione di un metodo o di una classe si ripercuote su più
    diagrammi contemporaneamente.
    \item \textbf{Programmatore}: erano state rendicontate solo le ore
    effettivamente dedicate alla produzione di codice, non quelle di
    progettazione preliminare per quanto riguarda script di automazione (es. script per pedici 
    glossario e calcolo dell'indice di gulpease), form di rendicontazione e pipeline;
     svolta prima dello sviluppo vero e proprio.
    \item \textbf{Verificatore}: non era stato rendicontato il tempo di lettura
    di un documento, che è parte integrante del processo di verifica, né,
    per il codice, il tempo di esecuzione dei test in locale; si
    contabilizzava solo il tempo strettamente legato alla compilazione
    della Pull Request, alla richiesta di cambiamenti e all'approvazione.
\end{itemize}

Il ruolo di Responsabile non compare in questo elenco e non presenta
alcuno scostamento: è risultato l'unico ruolo in cui la rendicontazione
era già corretta e allineata tra tutti i membri, poiché associata a
tipologie di task con una durata di riferimento chiara e condivisa fin
dall'inizio del progetto (ad esempio, la stesura di un verbale, circa 60
minuti, o l'aggiornamento del Diario di Bordo, circa 20 minuti). Per
questo motivo le ore rendicontate come Responsabile non sono state
modificate dalla rettifica.

\subsection{Risultati per ruolo}

Per ciascun ruolo interessato, la Tabella~\ref{tab:appendice-delta} riporta lo
scostamento complessivo (in ore), il numero di sprint a cui fanno riferimento le 
attività analizzate, il valore medio per sprint, e la motivazione individuata.

\begin{table}[H]
\centering
\renewcommand{\arraystretch}{1.25}
\footnotesize
\begin{tabular}{|p{2.2cm}|>{\centering\arraybackslash}p{1.2cm}|>{\centering\arraybackslash}p{1.3cm}|>{\centering\arraybackslash}p{1.6cm}|p{6.3cm}|}
\hline
\textbf{Ruolo} & \textbf{Ore} & \textbf{N. sprint\textsuperscript{*}} & \textbf{Media h/sprint} & \textbf{Motivazione} \\
\hline
Responsabile & 0{,}00 & --- & --- & Nessuno scostamento sistematico individuato. \\
\hline
Amministratore & 34{,}50 & 21 & 1{,}64 & Tempo non rendicontato per creare/raccogliere i dati dei grafici del PdQ e le tabelle dei consuntivi del PdP, e per la revisione/sistemazione di paragrafi ed errori nella gestione delle NdP. \\
\hline
Analista & 20{,}83 & 16 & 1{,}30 & Tempo non rendicontato per la creazione dei grafici dei casi d'uso e per la rinumerazione manuale dei casi d'uso ad ogni aggiunta/eliminazione. \\
\hline
Progettista & 16{,}33 & 9 & 1{,}81 & Tempo non rendicontato per la creazione/modifica di diagrammi delle classi, package, sequenza e deployment, impattati anche da una singola modifica di denominazione. \\
\hline
Programmatore & 13{,}00 & 8 & 1{,}63 & Rendicontata solo la produzione di codice, non la progettazione preliminare di script, form di rendicontazione e pipeline. \\
\hline
Verificatore & 57{,}58 & 21 & 2{,}74 & Non rendicontati il tempo di lettura dei documenti in verifica e, per il codice, l'esecuzione dei test in locale; contabilizzata solo l'attività diretta sulla PR. \\
\hline
\end{tabular}

\smallskip
\footnotesize\textsuperscript{*}\textit{Amministratore e Verificatore: tutti i 21 sprint del progetto. Analista: 16 sprint (8 RTB + 8 PB). Progettista: 9 sprint (tutti quelli della PB). Programmatore: 8 sprint (2 PoC + 3 MVP + 3 dedicati a CI/metriche/script).}

\caption{Scostamento del tempo per ruolo}
\label{tab:appendice-delta}
\end{table}

\subsection{Distribuzione per membro}

La Tabella~\ref{tab:appendice-delta-membro} riporta la stessa analisi
suddivisa per membro del gruppo, in ore.

\begin{table}[H]
\centering
\renewcommand{\arraystretch}{1.3}
\begin{tabularx}{\textwidth}{|X|c|c|c|c|c|c|}
\hline
\textbf{Ruolo} & \textbf{Armando} & \textbf{Filippo} & \textbf{Francesco} & \textbf{Ricardo} & \textbf{Roberto} & \textbf{Sofia} \\
\hline
Responsabile   & 0{,}00 & 0{,}00 & 0{,}00 & 0{,}00 & 0{,}00 & 0{,}00 \\
Amministratore & 6{,}25 & 2{,}50 & 8{,}50 & 6{,}17 & 11{,}08 & 0{,}00 \\
Analista       & 7{,}50 & 7{,}00 & 3{,}08 & 1{,}08 & 2{,}17 & 0{,}00 \\
Progettista    & 0{,}00 & 5{,}58 & 2{,}08 & 3{,}00 & 5{,}67 & 0{,}00 \\
Programmatore  & 2{,}00 & 3{,}50 & 1{,}00 & 5{,}00 & 1{,}50 & 0{,}00 \\
Verificatore   & 10{,}42 & 15{,}83 & 9{,}17 & 7{,}33 & 14{,}83 & 0{,}00 \\
\hline
\textbf{Totale} & \textbf{26{,}17} & \textbf{34{,}41} & \textbf{23{,}83} & \textbf{22{,}58} & \textbf{35{,}25} & \textbf{0{,}00} \\
\hline
\end{tabularx}
\caption{Scostamento del tempo per membro (ore)}
\label{tab:appendice-delta-membro}
\end{table}

Il totale complessivo dello scostamento individuato è di 142{,}25 ore, 
distribuito in modo non uniforme tra i membri in funzione
delle attività effettivamente svolte da ciascuno nei ruoli interessati.
Sofia non presenta scostamenti in nessun ruolo: come indicato in
Sezione~\ref{sec:appendice-rendicontazione}, ha valutato la propria
rendicontazione già coerente con il concetto di ore produttive maturato
dal gruppo, e ha scelto consapevolmente di non modificarla.

\subsection{Conclusioni}

Lo scostamento individuato in questa sede deriva da una non completa comprensione, 
da parte della maggior parte dei membri, di quali attività fossero effettivamente 
da rendicontare per il proprio ruolo. Il gruppo documenta questa
analisi, e la scelta di ciascun membro di riallineare o meno la propria
rendicontazione, a beneficio della trasparenza sul metodo adottato.